\section{Objetivos y contribuciones}\label{sec:conclusiones:evaluacion}

La Tabla~\ref{tab:conclusiones:evaluacion} da el estado de cada objetivo específico y de la contribución que le corresponde, con los epígrafes de la memoria que lo sostienen.
OBJ-03 y OBJ-04 se justifican primero, porque el Capítulo~\ref{cap:introduccion} sitúa en ellos el componente de investigación.

\begin{table}[tbp]
\centering
\small
\setlength{\tabcolsep}{3pt}
\begin{tabular}{|>{\raggedright\arraybackslash}p{3.0cm}|c|c|>{\raggedright\arraybackslash}p{2.0cm}|>{\raggedright\arraybackslash}p{2.2cm}|>{\raggedright\arraybackslash}p{3.0cm}|}
\hline
\textbf{Objetivo} & \textbf{RQ} & \textbf{Contrib.} & \textbf{Estado del objetivo} & \textbf{Estado de la contribución} & \textbf{Epígrafes que lo sostienen} \\
\hline
OBJ-01. Selección justificada del dominio & RQ1 & C-01 & cumplido en parte & cumplida en parte & \ref{sec:propuesta:seleccion}, \ref{sec:propuesta:respuesta} \\
\hline
OBJ-02. Construcción del catálogo de Datasheets & RQ2 a RQ4 & C-01 & cumplido & cumplida en parte & \ref{sec:propuesta:trazabilidad}, \ref{sec:propuesta:prime4api}, \ref{sec:esquema:contrato} \\
\hline
OBJ-03. Evolución del esquema & RQ2 & C-02 & cumplido en parte & cumplida & \ref{sec:esquema:inventario}, \ref{sec:esquema:fuera}, \ref{sec:esquema:respuesta} \\
\hline
OBJ-04. Análisis de la capacidad & RQ3 & C-03 & cumplido & cumplida & \ref{sec:limites:geometrica}, \ref{sec:limites:catalogo}, \ref{sec:limites:respuesta} \\
\hline
OBJ-05. Método de comparación & RQ4 & C-04 & cumplido & cumplida & \ref{sec:comparacion:catalogo}, \ref{sec:comparacion:procedimiento}, \ref{sec:comparacion:respuesta} \\
\hline
OBJ-06. Contribución a SPHERE & ninguna & C-01 & cumplido & cumplida en parte & \ref{sec:propuesta:sphere} \\
\hline
\end{tabular}
\caption[Estado de cada objetivo específico y de la contribución que le corresponde]{Estado de cada objetivo específico y de la contribución que le corresponde. RQ: la pregunta de investigación a la que sirve el objetivo. C-01: catálogo de Datasheets de APIs de correo. C-02: evolución de ACM2YAML. C-03: caracterización de la relación entre límites y \textit{rates}. C-04: método de comparación sistemática. OBJ-02 sirve de medio a RQ2, RQ3 y RQ4, porque el catálogo es la condición para responderlas. C-01 aparece tres veces porque responde a OBJ-01, OBJ-02 y OBJ-06. OBJ-06 no deriva de ninguna RQ.}
\label{tab:conclusiones:evaluacion}
\end{table}

\textbf{Evolución del esquema (OBJ-03 y C-02).}
De los 35 elementos no soportados que reúne el inventario de las ocho APIs, 24 superan el criterio de promoción: aparecen en dos APIs o más, o una ejecución de PRIME4API muestra su efecto sobre la Datasheet (Sección~\ref{sec:esquema:inventario}).
Hay nueve propuestas de extensión.
Ocho responden a elementos promovidos, como P-01, los seis volúmenes contratados que Postmark vende dentro de un mismo plan.
La novena, P-05, el plan con precio que existe y no es público, sale del fallo del motor ante el precio \texttt{custom} del plan Enterprise de Resend.

Los otros 16 elementos promovidos no tienen propuesta.
Doce solo tienen una forma esbozada que no se ha contrastado con la guía de ACM2YAML (Tabla~\ref{tab:sinforma}).
Los cuatro restantes quedan fuera de las propuestas por otras razones, como los \textit{add-ons} de IP dedicada, que el esquema deja al canal web (Sección~\ref{sec:esquema:fuera}).
OBJ-03 queda cumplido en parte, porque pide la extensión de cada carencia identificada y esos doce promovidos no la tienen fijada.

C-02 se cumple: cada una de las siete propuestas que amplían lo que ACM2YAML v0.3 puede expresar responde a un caso real que el esquema no representa, como el tope con que SendGrid limita el envío mensual a diez veces el volumen contratado, sobrecoste incluido.
Las otras dos quedan fuera del enunciado de C-02: P-04, la marca de valor no publicado junto al campo, y P-05.
ACM2YAML v0.3 admite esos valores, y es PRIME4API el que falla al leerlos.

\textbf{Análisis de la capacidad (OBJ-04 y C-03).}
La cuota es la restricción activa en Postmark, Resend, SendGrid, MailerSend, Brevo y Mailgun, seis de las ocho APIs (Sección~\ref{sec:limites:respuesta}).
La curva de capacidad $C(t)$ que dibuja PRIME4API (Figura~\ref{fig:capacidad}) da el máximo de correos que se pueden enviar desde el inicio de la ventana hasta cada instante sin pasar ningún límite.
Su codo es el instante en que deja de mandar el \textit{rate} y empieza a mandar la cuota, y con un destinatario por petición llega a los 83 minutos en Resend Pro y a los 41 en MailerSend Hobby.
Dos límites del catálogo no caben en esa curva: los dos contadores de Mailgun sobre la misma ventana y la franquicia sin ventana de Mailchimp Transactional, API en la que manda el caudal.

OBJ-04 y C-03 se cumplen sobre una base estrecha.
Solo en esas seis APIs la Datasheet da cifra al \textit{rate} y a la cuota de un mismo plan, y los valores medidos en MailerSend, Amazon SES y Mailchimp Transactional salen de cuentas sin aprobar o en pruebas, no del plan contratado.

\textbf{Método de comparación (OBJ-05 y C-04).}
El catálogo de preguntas reúne 32 preguntas de decisión, y el procedimiento que instancia cada una sobre cada API se ha aplicado a las ocho (Sección~\ref{sec:comparacion:procedimiento}).
OBJ-05 y C-04 se cumplen con tres salvedades.
La primera es que las cuatro preguntas que piden HARVEY4API cuentan como respondibles por la especificación de HARVEY4API que recoge Ruiz López \citep{ruizlopez2026tfm}, y solo se contrastaron con las respuestas que dio el modelo de lenguaje Claude Opus 5.5 al leer las Datasheets.

La segunda salvedad es que las diez respuestas de Mailchimp Transactional que resuelve una operación de PRIME4API están calculadas a mano.
El motor devuelve HTTP 500 en 14 llamadas por dos campos vacíos que el esquema admite: el tope de destinatarios y el periodo de la franquicia de la cuenta de prueba.
La tercera salvedad es que una de las tres preguntas funcionales, la retención de los registros de envío, pide un campo que ACM2YAML v0.3 no tiene.
Las otras dos se responden leyendo la matriz de \textit{features} por plan.

\textbf{Catálogo, selección y SPHERE (C-01 frente a OBJ-01, OBJ-02 y OBJ-06).}
OBJ-02, la construcción del catálogo, se cumple.
Los términos de consumo de cada proveedor quedan recogidos con su URL y su fecha de acceso, y cada valor de las ocho Datasheets lleva la cita literal que lo sostiene (Sección~\ref{sec:propuesta:trazabilidad}).
En Postmark no se archivaron las páginas que cita su Datasheet, y su versión se conserva solo como URL y fecha de acceso.

OBJ-01, la selección del dominio, se cumple en parte.
El protocolo parte de tres fuentes de inclusión: las APIs del trabajo de Noguerol Tirado \citep{noguerol2025tfm}, las ya modeladas en el ecosistema y una consulta a un modelo de lenguaje (Sección~\ref{sec:propuesta:fuentes}).
La primera no aportó ninguna API, porque se tuvo acceso a la memoria de Noguerol Tirado cuando la selección ya estaba cerrada.
La consulta presentaba Amazon SES como pago por uso puro, cuando ya publicaba planes, y repetirla no garantiza las mismas APIs (Sección~\ref{sec:propuesta:respuesta}).
Postmark, el piloto con el que se probó y se ajustó el método, no entra por ninguna de las tres fuentes.

OBJ-06, la contribución a SPHERE \citep{sphere}, se cumple.
Los ocho pricings en Pricing2YAML pasan el analizador de pricing4ts 0.10.4, la versión que usa la plataforma, y cada una de las 20 carencias del formato tiene su caso real.
Pricing2YAML no distingue un caudal de una cuota: en Resend Free, el caudal de 10 peticiones por segundo y la cuota de 3.000 correos al mes se escriben igual, como un límite renovable con su periodo, y la analítica de SPHERE los cuenta igual (Sección~\ref{sec:propuesta:sphere}).
La maqueta de la colección en SPHERE propone cómo mostrar los términos de consumo en la ficha de cada API.\footnote{La maqueta puede abrirse en \url{https://danvelcam.github.io/tfm-pricing-apis-correo/sphere/coleccion.html} (Sección~\ref{sec:propuesta:sphere:coleccion}).}

C-01 se cumple en parte por tres razones.
La primera es que la colección está cargada en SPHERE desde el 27 de octubre de 2026, pero es privada hasta el depósito.
La segunda es que ninguna Datasheet supera todas las comprobaciones que se le hacen al ingerirla en PRIME4API (Tabla~\ref{tab:verificacion}).
La tercera es que HARVEY4API, la otra herramienta del ecosistema que debería reutilizar las Datasheets, no se ha ejecutado.

En cinco de las ocho Datasheets, las comprobaciones fallan por un valor ausente o por un precio no público.
En Postmark queda vacío el tope de destinatarios, que el proveedor no pone, y en Mailchimp Transactional ese mismo tope y el periodo de la franquicia de la cuenta de prueba.
Fallan diez comprobaciones en cada una, y cinco en Resend, que no publica el \textit{rate} de su plan Enterprise.
Un precio que existe y no es público invalida trece comprobaciones en SendGrid, trece en Brevo y ocho en Resend.
MailerSend, Amazon SES y Mailgun solo fallan la comprobación de la curva de capacidad acumulada, con el mismo mensaje de error.
En MailerSend y Amazon SES el fallo es del motor, y en Mailgun no se ha aislado su causa.
